\begin{figure}[ht]
\centering
\begin{tikzpicture}[x=1cm, y=1cm, font=\small]

% Thermistor resistance-temperature curve R(T) = R0 exp(B(1/T - 1/T0)),
% strongly convex and falling. Show the curve, an operating point, and
% the tangent (small-signal) line through it. Illustrative shape only.
% T axis (deg C) from 0 to 80 mapped to x in [0, 9]; x = 9*(T)/80.
% R axis from 0 to 30 kOhm mapped to y in [0, 5]; y = 5*R/30.
% Curve sampled at chosen (T, R) points to show convex falling shape.

\draw[->, thick] (-0.3, 0) -- (9.6, 0)
  node[right] {$T\;(^{\circ}\mathrm{C})$};
\draw[->, thick] (0, -0.3) -- (0, 5.6)
  node[above] {$R\;(\mathrm{k}\Omega)$};

% x ticks
\foreach \T in {0,20,40,60,80} {
  \pgfmathsetmacro\xx{9*\T/80}
  \draw (\xx, 0) -- (\xx, -0.12);
  \node[anchor=north, font=\scriptsize] at (\xx, -0.12) {\T};
}
% y ticks
\foreach \R in {0,10,20,30} {
  \pgfmathsetmacro\yy{5*\R/30}
  \draw (0, \yy) -- (-0.12, \yy);
  \node[anchor=east, font=\scriptsize] at (-0.14, \yy) {\R};
}

% Curve: representative NTC points (T degC, R kOhm)
% (0,28) (10,20) (20,14.4) (25,12.0) (30,10.2) (40,7.4) (50,5.4)
% (60,4.0) (70,3.0) (80,2.3). Plot scaled.
\draw[very thick, engblue]
  plot coordinates {
    ({9*0/80},  {5*28.0/30})
    ({9*10/80}, {5*20.0/30})
    ({9*20/80}, {5*14.4/30})
    ({9*25/80}, {5*12.0/30})
    ({9*30/80}, {5*10.2/30})
    ({9*40/80}, {5*7.4/30})
    ({9*50/80}, {5*5.4/30})
    ({9*60/80}, {5*4.0/30})
    ({9*70/80}, {5*3.0/30})
    ({9*80/80}, {5*2.3/30})
  };

% Operating point at T0 = 25 degC, R0 = 12.0 kOhm.
% Local slope dR/dT approx (10.2 - 14.4)/(30-20) = -0.42 kOhm/degC.
% Tangent line over T in [10, 45]: R = 12.0 - 0.42*(T - 25).
% At T=10: R = 12.0 - 0.42*(-15) = 18.3; at T=45: R = 12.0 - 0.42*20 = 3.6.
\draw[thick, engred, dashed]
  plot coordinates {
    ({9*10/80}, {5*18.3/30})
    ({9*45/80}, {5*3.6/30})
  };

\fill[engblack] ({9*25/80}, {5*12.0/30}) circle (1.8pt);
\node[engblack, anchor=south west, font=\scriptsize]
  at ({9*25/80 + 0.1}, {5*12.0/30 + 0.05})
  {operating point $(T_{0}, R_{0})$};

% small-signal box near operating point
\draw[thin, enggray, densely dotted]
  ({9*20/80}, 0) -- ({9*20/80}, {5*14.4/30});
\draw[thin, enggray, densely dotted]
  ({9*30/80}, 0) -- ({9*30/80}, {5*10.2/30});
\node[enggray, anchor=north, font=\tiny]
  at ({9*25/80}, -0.45) {small-signal band};

\node[engblue, anchor=east, font=\scriptsize] at (9.4, 4.7)
  {$R(T)$ (NTC)};
\node[engred, anchor=east, font=\scriptsize] at (9.4, 4.3)
  {tangent at $T_{0}$};

\end{tikzpicture}
\caption[Operating-point linearisation of a thermistor]{The
resistance-temperature curve of a negative-temperature-coefficient
thermistor (blue) falls steeply and convexly with temperature. A
measurement circuit operated in a narrow band around the operating
point $(T_{0}, R_{0})$ replaces the curve by its tangent (red dashed),
the small-signal model $\delta R \approx (\mathrm{d}R/\mathrm{d}T)_{T_{0}}
\,\delta T$. The slope at the operating point sets the sensor's local
sensitivity in $\mathrm{k}\Omega$ per degree; the curvature, the
second derivative, sets the width of the band over which the linear
model stays inside a stated tolerance. Reading the design off the
derivative at one point is the entire content of small-signal sensor
calibration.}
\label{fig:vol02:ch05:operating-point}
\end{figure}
